\documentclass[12pt]{article}
\usepackage[margin=1in]{geometry}
\usepackage{amsmath}
\usepackage{amssymb}
\usepackage{amsthm}
\usepackage{enumitem}
\usepackage{xcolor}
\usepackage{pagecolor}
\usepackage{marginnote}
\usepackage{fancyhdr}
\usepackage{bm}


% Define cream background color
\definecolor{cream}{HTML}{faf7f1}
\pagecolor{cream}

% Define red color for dates
\newcommand{\red}[1]{\textcolor{red}{#1}}

% Define defn and defeq commands
\newcommand{\defn}[1]{\textbf{#1}}
\newcommand{\defeq}{\mathrel{\vcenter{\baselineskip0.5ex \lineskiplimit0pt \hbox{\scriptsize.}\hbox{\scriptsize.}}}=}

% Define sidenote command (simplified - uses marginnote)
\newcommand{\sidenote}[1]{\marginnote{\small\itshape #1}}

% Title formatting
\makeatletter
\renewcommand{\maketitle}{%
  \begin{flushleft}
    {\Large\bfseries CS173: Discrete Structures}\\[0.3em]
    {\Large\bfseries Problem Set 0}\\[0.5em]
    {\normalsize Mohammed Al Salhi}
  \end{flushleft}
  \vspace{1em}
}
\makeatother
\AtBeginDocument{\mathversion{bold}}

\begin{document}

\maketitle

\begin{enumerate}[left=0pt, itemsep=2em, parsep=1em]
  \item
    Determine the truth value, if any, of each of the following sentences.
    
    \begin{enumerate}[label=(\alph*), itemsep=1.5em, leftmargin=2em]
      \item
        ``Every integer is either even or odd.''
        
        \textbf{TRUE.}
        
        Let $n$ be an integer. Thus $2n$ is an even number, and $2n+1$ is and an odd number. Incrementing up from $2n+1$ is $2n+2 = 2k$ where $k=n+1$ and is an integer meaning $2k$ is again even. By the principles of induction all integer are thus even or odd. 
    
      \item
        ``Every number is either even or odd.''
        
        \textbf{FALSE.}
        
        Can be proven by counterexample. $2.5$ is a number and is neither even nor odd. 
        
      \item
        ``Every non-empty set of integers has a smallest element.''
        
        \textbf{FALSE.}
        
        Can be proven by counterexample. Take the set of all negative integers $\{\ldots, -2, -1\}$. This set has no smallest element.
        
      \item
        ``Every non-empty set of natural numbers has a smallest element.''
        
        \textbf{TRUE.}
        
        The set of natural numbers is bounded below by $0$ hence all subsets of natural numbers must have a smallest value. 
        
      \item
        ``Every natural number is boring.''
        
        \textbf{FALSE.}
        
        The premise of being boring is dialectical. Fore there to be 'boring' natural numbers there has to an opposite, or a reference. That suggest that there must exist some 'interesting' or 'non-boring' natural numbers for 'boring' natural numbers to exist. If all numbers were boring, then none would be boring. The term would lose meaning.
        
      \item
        ``No natural number is boring.''
        
        \textbf{FALSE.}
        
        From part (e), there has to exist some boring natural numbers so that non-boring' natural numbers exist.
    \end{enumerate}

  \item
    Determine the truth value, if any, of each of the following sentences.
    
    \begin{enumerate}[label=(\alph*), itemsep=1.5em, leftmargin=2em]
      \item
        ``This sentence is true.''
        
        \textbf{INDETERMINATE.}
        
        Let $S \defeq$ ``$S$ is true.'' 
        If $S$ is true, then the assertion is true, which is consistent.
        If $S$ is false, then what $S$ assertion is false, which is also consistent.
        
      \item
        ``This sentence is false.''
        
        \textbf{PARADOX.}
        
        Let $Q \defeq$ ``$Q$ is false.'' If $Q$ is true, then what $Q$ asserts must hold, meaning $Q$ is false which is a contradiction.
        If $Q$ is false, then what $Q$ asserts does not hold, meaning $Q$ is not false, so $Q$ is true which again is a contradiction.
        
      \item
        ``The set of all sets that contain themselves contains itself.''
        
        \textbf{INDETERMINATE.}
        

        Let this set be $R$. If $R \in R$, then $R$ satisfies the definition of $R$ and thus is in $R$, which is consistent.
        If $R \notin R$, then $R$ does not satisfy the definition, which is also consistent.
        
      \item
        ``The set of all sets that do not contain themselves does not contain itself.''
        
        \textbf{PARADOX.}
        
        Let this set be $N$. If $N \notin N$, then $N$ satisfies the definition of $N$ and thus is in $N$, which is a contradiction.
        If $N \in N$, then $N$ does not satisfy the definition, but is in $N$ which is also a contradiction.

        
      \item
        ``If this sentence is false, then $7$ is a prime number.''
        
        \textbf{TRUE.}
        
        Let $A \defeq$ ``If $A$ is false, then 7 is a prime number.'' and $B \defeq$ ``7 is prime''. Then $A \defeq$ $\neg A \Rightarrow B$.
        If $A$ is false then $\neg A$ is true, which implies $B$ has to be true. $B$ is obviously true, which means this is consistent. 
        If $A$ is true then $\neg A$ is false, which means $B$ is not dependent on $A$ and can be either $\top$ or $\bot$, thus this is still consistent. 
        
        
      \item
        ``If this sentence is true, then $7$ is not a prime number.''
        
        \textbf{FALSE.}
        
        Let $A \defeq$ ``If $A$ is true, then $7$ is not a prime number.'' and $B \defeq$ ``$7$ is not prime''. Then $A \defeq A \Rightarrow B$.
        If $A$ is true then it implies $B$ has to be true. But $B$ is obviously false, so this is inconsistent.
        If $A$ is false then it means $B$ is not dependent on $A$ and can be either $\top$ or $\bot$, thus this is consistent. So $A$ has to be $\bot$.

    \end{enumerate}

  \item
    Consider the infinite sequence of sentences indexed by natural numbers
    \begin{equation*}
      S_0, S_1, S_2, \ldots, S_i, \ldots
    \end{equation*}
    where each sentence asserts that all of the subsequent sentences are false.
    \begin{equation*}
      S_i \defeq \text{``}S_j \text{ is false for all } j > i\text{''}
    \end{equation*}
    For example, $S_3 \defeq$ ``$S_j$ is false for all $j > 3$,'' which means $S_3$ is equivalent to the assertion that ``the sentences $S_4, S_5, S_6, \dots$ are all false.''
    
    What are the truth values, if any, of the sentences in this sequence?
    

    \textbf{All $S_n$ are INDETERMINATE.}
    
    $S_0$ cannot be $\top$, because if $S_0$ were true, then it would assert that $S_1$ is $\bot$.  
    But $S_1$ says that all sentences $S_n$ with $n>1$ are false.  
    If $S_1$ is false, then there exists some $S_n$ with $n>1$ that is true, which contradicts the assertion made by $S_0$ that all subsequent sentences are false.
    
    The same applies to any $S_n$: assuming $S_n$ is $\top$ forces $S_{n+1}$ to be $\bot$, which in turn implies the existence of some later sentence that is $\top$, contradicting the original claim of $S_n$.
    
    The only way all sentences could avoid contradiction by being true would be if there were a final sentence $S_N$ that is true. However, there is no last sentence, since the sequence is infinite.
    
    Therefore, none of the sentences $S_n$ can be assigned a truth value.
\end{enumerate}

\end{document}